\section{Index a complete file before choosing a frame}
A record routine sees one frame. Opening an animation also requires checks that
span frames: the global header, the declared frame count, the sequence of record
windows and the exact file end. Implement \texttt{video\_index} for this task.
For now its caller supplies the complete file as bytes in memory. The later
loader must obtain those bytes from a file and establish the actual end of input;
this function performs no file-system calls.

Pass the file address in \texttt{\$a0}, its complete byte count in
\texttt{\$a1}, a native offset-directory address in \texttt{\$a2}, and an
80-byte native context in \texttt{\$a3}. The file address may be unaligned.
The directory and context are eight-byte aligned. A directory entry is an
unsigned eight-byte offset from the file beginning, not a pointer. Up to 256
entries require 2048 bytes. A later seek can obtain a record's address by adding
its admitted offset to the retained file base.

The context supplies a separate scratch image. The illustrated validator calls
\texttt{video\_apply\_record} on each record using this image. A full first
record initializes it; each following record then has a valid previous image.
The scratch image is temporary work storage. Its active contents are unspecified
after either success or failure and must not be used as the displayed frame.
This permits rejection of a late record after earlier records have changed the
scratch image. Public outputs still remain unchanged on failure.

\begin{table}[htbp]
  \centering
  \caption{Index context. All fields use eight native bytes.}
  \label{tab:video-index-context}
  \begin{tabular}{r p{1.5in} p{2.8in}}
    \toprule
    Offset & Field                & Meaning                                                    \\
    \midrule
    0      & Directory capacity   & Input byte capacity, at least eight times the frame count. \\
    8      & Scratch address      & Input pointer to four-byte-aligned image storage.          \\
    16     & Scratch capacity     & Input byte capacity, at least 16\,384.                     \\
    24     & Frame count          & Output, one through 256.                                   \\
    32     & Full-frame period    & Output, one through 64.                                    \\
    40     & Duration             & Output milliseconds per frame, one through 1000.           \\
    48     & Validated file bytes & Output, the exact admitted input count.                    \\
    56     & Directory entries    & Output, equal to the frame count.                          \\
    64--79 & Retained suffix      & Leave both eight-byte fields unchanged.                    \\
    \bottomrule
  \end{tabular}
\end{table}

Success returns zero in \texttt{\$v0} and the active directory byte count in
\texttt{\$v1}. Failure returns one and zero, retaining the complete directory
and context, including earlier successful outputs. Keep the file immutable on
both paths. Scratch accesses must fit its 16\,384 active bytes. Retain every
other caller byte, saved register, frame pointer and stack pointer.

The caller promises readable file bytes, sufficient writable directory storage
for its declared capacity, a writable context, scratch storage and a valid stack.
These windows are disjoint. Use a sixteen-byte-aligned stack at entry and keep
that alignment across nested calls. The supplied caller clears the low four
stack-address bits before its first call. The illustrated index frame is 256
bytes; it saves the return address because calling another function replaces
\texttt{\$ra}. Its 48-byte record context begins at stack offset 128 and stays
separate from the record routine's own frame below the current stack pointer.

Build this stage from your completed record working file. Run
\texttt{python3 packages/034-lossless-video/construction/video-index-construction/attach.py --source work.asm
  --output index-work.asm}. The attachment retains your record instruction text
and gives its old example entry the name \texttt{record\_example\_main}; the
new \texttt{main} calls the indexer. It refuses conflicting labels and existing
output files. Prepare an observation with
\texttt{python3 packages/034-lossless-video/construction/video-index-construction/prepare.py s24-1 --source
  index-work.asm --output observe.asm}. Preparation changes data, preserving
instructions. An unfinished record prerequisite propagates development status
four through the indexer. Complete that prerequisite before interpreting the
later checkpoints.

Use two passes over the retained file. The first validates every record and exact
completion, publishing nothing. The second records the already-validated offsets.
The file must remain immutable between the passes and while the directory is in
use. A validation-only implementation may avoid scratch reconstruction if it
establishes all the same format checks. Free-writing grading checks the public
behavior and scratch bounds, without prescribing a private frame or scratch image
contents. Development status four remains limited to unfinished construction;
after the directory-writing region is filled, an unfinished later region may
return four with a partial directory. A completed failure status one never does.
